% !TEX program = xelatex
\documentclass{booklet_style}

%\usepackage[pagewise, mathlines]{lineno}

%\usepackage[draft]{todonotes}
% \usepackage[disable]{todonotes} % notes not showed

%\setmainfont{Garamond}

\usepackage[hyperref=true,
            backref=true,
            style=numeric-comp,
            sorting=none,
            maxnames=100,
            giveninits=true,
            url=false,
            isbn=false,
            defernumbers=true
            ]{biblatex}

\renewbibmacro{in:}{}
\AtEveryBibitem{%
    \clearlist{language}%
}

\addbibresource{sajat.bib}
\addbibresource{EMBL.bib}

\sloppy
\frenchspacing

\DeclareRefcontext{journal}{labelprefix=J}
\DeclareBibliographyCategory{journal}
\addtocategory{journal}{balazs_real-time_2017,de_medeiros_light-sheet_2016,strnad_inverted_2016,hoyer_breaking_2016}


\DeclareRefcontext{conference}{labelprefix=C}
\DeclareBibliographyCategory{conference}
\addtocategory{conference}{balazs_gpu-based_2016,balazs_gpu-based_2016-1,balazs_gpu-based_2017}

\def\b3d{B\textsuperscript{3}D}

\author{\href{mailto:balint.balazs@embl.de}{Bálint Balázs}}
\supervisor{\href{mailto:hufnagel@embl.de}{Lars Hufnagel}, \href{mailto:rozsabal@koki.hu}{Balázs Rózsa}}
\lab{\href{http://www.embl.de/}{European Molecular Biology Laboratory}}
\university{\href{http://www.ppke.hu/}{Pázmány Péter Catholic University}}
\collegeordept{\href{http://www.itk.ppke.hu/}{Faculty of Information Technology}}
\degreedate{\textit{Theses of the Ph.D. Dissertation} \\ Budapest, 2017}
\ppke{\includegraphics[width=1cm]{figures/0_front/ITK_logo}}
\embl{\includegraphics[width=2.5cm]{figures/0_front/EMBL_logo}}
\degree{PhD}

\title{Title goes here}
%\date{2012 november}



\begin{document}
%\begin{linenumbers}
\maketitle
%\pagenumbering{roman}

\newpage
\setcounter{page}{1}
\section{Introduction}


\section{DualMouse SPIM}

\section{\b3d image compression algorithm}
other papers cited here
huisken \cite{huisken_optical_2004}
muvi spim \cite{krzic_multiview_2012}

\section{New scientific results}

\paragraph{Thesis I.}\textit{I have designed and constructed a new light-sheet microscope suitable for high sensitivity imaging of delicate samples. This setup greatly improves existing solutions by employing two high NA objectives}

Corresponding publications: \cite{de_medeiros_light-sheet_2016},\cite{strnad_inverted_2016}, \cite{hoyer_breaking_2016}

More description coming. %\todo{add more stuff here}

\paragraph{Thesis II.} \textit{I have developed a GPU-based image procesing pipeline for multi-view light-sheet microscopy that enables real time fusion of opposing views.}

Corresponding publications: \cite{balazs_gpu-based_2016}, \cite{balazs_gpu-based_2016-1}, \cite{balazs_gpu-based_2017}

\subsection*{Thesis Group II.}

\paragraph{Thesis III.} \textit{I have  developed a new image compression algorithm that enables noise dependent lossy compression of light microscopy images, and can reach a compression ratio of 100 fold while preserving the results of downstream data analysis steps.}

Corresponding publications: \cite{balazs_real-time_2017}, \cite{balazs_gpu-based_2016}, \cite{balazs_gpu-based_2016-1}, \cite{balazs_gpu-based_2017}\\[1ex]
\indent
All microscopy images contain inherent noise, that is 

Trying some formulas inline $q=1\sigma$. And more formuals in their own line:
\begin{equation}
    x = \frac{-b \pm \sqrt{b^2 - 4ac}}{2a}
\end{equation}

\paragraph{Thesis IV.} \textit{I have  developed a new GPU-based image compression library, \b3d, that implements the algorithm described in Thesis III, and allows for real time compression of microscopy images with a throughput of up to 1 GB/s.}

Corresponding publications: \cite{balazs_real-time_2017}, \cite{balazs_gpu-based_2016}, \cite{balazs_gpu-based_2016-1}, \cite{balazs_gpu-based_2017}

\b3d is an efficient, GPU-based image compression library allowing lossless and noise dependent lossy compression of microscopy images. Since many high-speed microscopy methods generate immense amounts of data, easily reaching terabytes per expriment, image compression is especially important to efficiently deal with such datasets.


\section{Application of the results}

%\bibliographystyle{IEEE/IEEEtran}
%\bibliography{sajat}
\newrefcontext{journal}
\printbibliography[category=journal, title={The Author's Journal Publications}]

\newrefcontext{conference}
\printbibliography[category=conference, title={The Author's Conference Presentations}]

\newrefcontext
\printbibliography[notcategory=journal,notcategory=conference, resetnumbers=true, title={References}]

%\end{linenumbers}
\end{document}